\documentclass[11pt,a4paper]{article}
\usepackage[T1]{fontenc}
\usepackage[utf8]{inputenc}
\usepackage{lmodern}
\usepackage[czech]{babel}
\usepackage[margin=2cm]{geometry}
\usepackage{amsmath,amssymb,mathtools}
\usepackage{xcolor,cancel,enumitem}
\usepackage{tikz,pgfplots}
\pgfplotsset{compat=1.18}
\usetikzlibrary{arrows.meta,calc,patterns,babel}
\pagestyle{empty}
\setlength{\parindent}{0pt}
\setlength{\parskip}{4pt}
\newcommand{\vysl}[1]{\boxed{\color{red!75!black}#1}}
\newcommand{\pozn}[1]{{\color{red!75!black}#1}}
\newcommand{\dx}{\,\mathrm{d}x}
\newcommand{\dt}{\,\mathrm{d}t}
\newcommand{\subst}[1]{\left|\begin{array}{l}#1\end{array}\right|}
\DeclareMathOperator{\arctg}{arctg}
\DeclareMathOperator{\arccotg}{arccotg}
\DeclareMathOperator{\tg}{tg}
\DeclareMathOperator{\cotg}{cotg}
\usepackage{tcolorbox}
\newtcolorbox{pomocny}{colframe=red!75!black, colback=white, coltitle=white, fonttitle=\bfseries, title={Pomocný výpočet:}, sharp corners}
\newcommand{\dy}{\,\mathrm{d}y}
% znaménkové schéma: \znam{bod/popis/typ,...}{znaménka v intervalech}
% typ: o = vyloučený bod (kroužek), x = obyčejný bod
\newcommand{\znam}[2]{%
\begin{tikzpicture}[baseline=-3pt]
\def\n{0}
\foreach \p/\l/\t [count=\i] in {#1} {\xdef\n{\i}}
\pgfmathsetmacro{\L}{1.6*(\n+1)}
\draw (0,0) -- (\L,0);
\foreach \p/\l/\t [count=\i] in {#1} {
  \pgfmathsetmacro{\xx}{1.6*\i}
  \ifnum\pdfstrcmp{\t}{o}=0 \draw[fill=white] (\xx,0) circle (2.5pt); \draw (\xx-0.09,0) -- (\xx+0.09,0) (\xx,-0.09) -- (\xx,0.09);
  \else \draw (\xx,0.1) -- (\xx,-0.1); \fi
  \node[below=2pt] at (\xx,0) {\small $\l$};
}
\foreach \s [count=\i] in {#2} {
  \pgfmathsetmacro{\xx}{1.6*\i-0.8}
  \node[above] at (\xx,0) {$\s$};
}
\end{tikzpicture}}
\newcommand{\aug}[1]{\left(\begin{array}{rrr|rrr}#1\end{array}\right)}
\newcommand{\krok}[1]{\text{\small\color{orange!80!black}#1}}
\renewcommand{\CancelColor}{\color{red!75!black}}

\begin{document}
\section*{DERIVACE FUNKCE}

\newcommand{\lnf}{\ln\frac{x+1}{x+2}}

\textbf{\textcircled{1}}\quad $y=e^{\sqrt{x^2+x+1}}$
\begin{align*}
y' &= e^{\sqrt{x^2+x+1}}\cdot\left(\sqrt{x^2+x+1}\right)' = e^{\sqrt{x^2+x+1}}\cdot\left(\left(x^2+x+1\right)^{\frac12}\right)' =\\
&= e^{\sqrt{x^2+x+1}}\cdot\frac12\left(x^2+x+1\right)^{\frac12-1}\cdot\left(x^2+x+1\right)' =\\
&= e^{\sqrt{x^2+x+1}}\cdot\frac12\left(x^2+x+1\right)^{-\frac12}\cdot(2x+1) = \vysl{\frac{(2x+1)\cdot e^{\sqrt{x^2+x+1}}}{2\sqrt{x^2+x+1}}}
\end{align*}

\hrule
\medskip
\textbf{\textcircled{2}}\quad $y=\sqrt{\arctg\lnf}\quad\left(=\left(\arctg\lnf\right)^{\frac12}\right)$
\begin{align*}
y' &= \frac12\left(\arctg\lnf\right)^{\frac12-1}\cdot\left(\arctg\lnf\right)' =\\
&= \frac12\left(\arctg\lnf\right)^{-\frac12}\cdot\frac{1}{1+\left(\lnf\right)^2}\cdot\left(\lnf\right)' =\\
&= \frac{1}{2\sqrt{\arctg\lnf}}\cdot\frac{1}{1+\ln^2\frac{x+1}{x+2}}\cdot\frac{1}{\frac{x+1}{x+2}}\cdot\left(\frac{x+1}{x+2}\right)' =\\
&= \frac{1}{2\sqrt{\arctg\lnf}}\cdot\frac{1}{1+\ln^2\frac{x+1}{x+2}}\cdot\frac{x+2}{x+1}\cdot\frac{(x+1)'\cdot(x+2)-(x+1)\cdot(x+2)'}{(x+2)^2} =\\
&= \frac{1}{2\sqrt{\arctg\lnf}}\cdot\frac{1}{1+\ln^2\frac{x+1}{x+2}}\cdot\frac{\cancel{x+2}}{x+1}\cdot\frac{1\cdot(x+2)-(x+1)\cdot1}{(x+2)^{\cancel{2}}} =\\
&= \frac{1}{2\sqrt{\arctg\lnf}}\cdot\frac{1}{1+\ln^2\frac{x+1}{x+2}}\cdot\frac{1}{x+1}\cdot\frac{\overbrace{x+2-x-1}^{1}}{x+2} =\\
&= \vysl{\frac{1}{2\sqrt{\arctg\lnf}\cdot\left(1+\ln^2\frac{x+1}{x+2}\right)\cdot(x+1)(x+2)}}
\end{align*}

\textbf{\textcircled{3}}\quad $y=\ln\dfrac{1}{x+\sqrt{x^2-1}}$
\begin{align*}
y' &= \frac{1}{\frac{1}{x+\sqrt{x^2-1}}}\cdot\Big(\underbrace{\frac{1}{x+\sqrt{x^2-1}}}_{\left(x+\sqrt{x^2-1}\right)^{-1}}\Big)' =\\
&= \left(x+\sqrt{x^2-1}\right)\cdot(-1)\left(x+\sqrt{x^2-1}\right)^{-1-1}\cdot\Big(x+\underbrace{\sqrt{x^2-1}}_{\left(x^2-1\right)^{\frac12}}\Big)' =\\
&= \left(x+\sqrt{x^2-1}\right)\cdot(-1)\left(x+\sqrt{x^2-1}\right)^{-2}\cdot\left(1+\frac12\left(x^2-1\right)^{\frac12-1}\cdot\left(x^2-1\right)'\right) =\\
&= \frac{\cancel{x+\sqrt{x^2-1}}}{-1\cdot\left(x+\sqrt{x^2-1}\right)^{\cancel{2}}}\cdot\left(1+\frac12\left(x^2-1\right)^{-\frac12}\cdot2x\right) =\\
&= \frac{-1}{x+\sqrt{x^2-1}}\cdot\left(1+\frac{x}{\sqrt{x^2-1}}\right) = \frac{-1}{\cancel{x+\sqrt{x^2-1}}}\cdot\frac{\cancel{\sqrt{x^2-1}+x}}{\sqrt{x^2-1}}
= \vysl{\frac{-1}{\sqrt{x^2-1}}}
\end{align*}

\hrule
\medskip
\textbf{\textcircled{4}}\quad $y=\ln\arctg\dfrac{\sin x}{1+\cos x}$
\begin{align*}
y' &= \frac{1}{\arctg\frac{\sin x}{1+\cos x}}\cdot\left(\arctg\frac{\sin x}{1+\cos x}\right)' =
\frac{1}{\arctg\frac{\sin x}{1+\cos x}}\cdot\frac{1}{1+\left(\frac{\sin x}{1+\cos x}\right)^2}\cdot\left(\frac{\sin x}{1+\cos x}\right)' =\\
&= \frac{1}{\arctg\frac{\sin x}{1+\cos x}}\cdot\frac{1}{1+\frac{\sin^2x}{(1+\cos x)^2}}\cdot\frac{(\sin x)'(1+\cos x)-\sin x(1+\cos x)'}{(1+\cos x)^2} =\\
&= \frac{1}{\arctg\frac{\sin x}{1+\cos x}}\cdot\frac{1}{\frac{(1+\cos x)^2+\sin^2x}{(1+\cos x)^2}}\cdot\frac{\cos x(1+\cos x)-\sin x(-\sin x)}{(1+\cos x)^2} =\\
&= \frac{1}{\arctg\frac{\sin x}{1+\cos x}}\cdot\frac{\cos x+\overbrace{\cos^2x+\sin^2x}^{1}}{1+2\cos x+\underbrace{\cos^2x+\sin^2x}_{1}} =\\
&= \frac{1}{\arctg\frac{\sin x}{1+\cos x}}\cdot\frac{\cancel{\cos x+1}}{\underbrace{2+2\cos x}_{2(\cancel{1+\cos x})}}
= \vysl{\frac{1}{2\arctg\frac{\sin x}{1+\cos x}}}
\end{align*}

\hrule
\medskip
\textbf{\textcircled{5}}\quad $y=\dfrac{1}{\log\left(3x^2+x+1\right)}=\left(\log\left(3x^2+x+1\right)\right)^{-1}$
\begin{align*}
y' &= (-1)\cdot\left(\log\left(3x^2+x+1\right)\right)^{-2}\cdot\left(\log\left(3x^2+x+1\right)\right)' =\\
&= \frac{-1}{\log^2\left(3x^2+x+1\right)}\cdot\frac{1}{\left(3x^2+x+1\right)\cdot\ln10}\cdot\left(3x^2+x+1\right)' =\\
&= \vysl{\frac{-(6x+1)}{\left(3x^2+x+1\right)\cdot\ln10\cdot\log^2\left(3x^2+x+1\right)}}
\end{align*}

\hrule
\medskip
\textbf{\textcircled{6}}\quad $y=5^{x^2-2x+1}$
\[
y' = 5^{x^2-2x+1}\cdot\ln5\cdot\left(x^2-2x+1\right)' = \vysl{(2x-2)\ln5\cdot5^{x^2-2x+1}}
\]

\hrule
\medskip
\textbf{\textcircled{7}}\quad $y=\dfrac{\cos x}{2\sin^2x}$
\begin{align*}
y' &= \frac{(\cos x)'\cdot2\sin^2x-\cos x\left(2\sin^2x\right)'}{\left(2\sin^2x\right)^2}
= \frac{-\sin x\cdot2\sin^2x-\cos x\cdot2\cdot2(\sin x)^{2-1}\cdot(\sin x)'}{4\sin^4x} =\\
&= \frac{-2\sin^3x-4\cos x\cdot\sin x\cdot(\cos x)}{4\sin^4x}
= \frac{-2\sin^3x-4\cos^2x\sin x}{4\sin^4x} =\\
&= \frac{-\cancel{2}\,\cancel{\sin x}\big(\overbrace{\sin^2x+\cos^2x}^{1}+\cos^2x\big)}{\underset{2}{\cancel{4}}\sin^{\cancel{4}\,3}x} = \vysl{-\frac{1+\cos^2x}{2\sin^3x}}
\end{align*}

\textbf{\textcircled{8}}\quad $y=\dfrac{2x^3}{\sqrt x}+\dfrac{x}{2\sqrt[3]x} = 2\,\dfrac{x^3}{x^{\frac12}}+\dfrac12\,\dfrac{x}{x^{\frac13}} = 2x^{\frac52}+\dfrac12x^{\frac23}$
\quad {\small (jen úpravy mocnin, pak až derivovat)}
\[
y' = 2\cdot\frac52x^{\frac52-1}+\frac12\cdot\frac23x^{\frac23-1} = 5x^{\frac32}+\frac13x^{-\frac13} = \vysl{5\sqrt{x^3}+\frac{1}{3\sqrt[3]x}}
\]

\hrule
\medskip
\textbf{\textcircled{9}}\quad $y=\dfrac{x}{\sqrt{1-x^2}}\cdot\arcsin x$
\begin{align*}
y' &= \left(\frac{x}{\sqrt{1-x^2}}\right)'\cdot\arcsin x+\frac{x}{\sqrt{1-x^2}}\cdot(\arcsin x)' =\\
&= \frac{(x)'\cdot\sqrt{1-x^2}-x\cdot\left(\sqrt{1-x^2}\right)'}{\left(\sqrt{1-x^2}\right)^2}\cdot\arcsin x+\frac{x}{\sqrt{1-x^2}}\cdot\frac{1}{\sqrt{1-x^2}} =\\
&= \frac{1\cdot\sqrt{1-x^2}-x\cdot\frac12\left(1-x^2\right)^{-\frac12}\cdot\left(1-x^2\right)'}{1-x^2}\cdot\arcsin x+\frac{x}{1-x^2} =\\
&= \frac{\sqrt{1-x^2}-\frac{x}{2\sqrt{1-x^2}}\cdot(-2x)}{1-x^2}\cdot\arcsin x+\frac{x}{1-x^2}
= \frac{\sqrt{1-x^2}+\frac{2x^2}{2\sqrt{1-x^2}}}{1-x^2}\cdot\arcsin x+\frac{x}{1-x^2} =\\
&= \frac{\dfrac{1\cancel{-x^2}\cancel{+x^2}}{\sqrt{1-x^2}}}{1-x^2}\cdot\arcsin x+\frac{x}{1-x^2}
= \vysl{\frac{1}{\sqrt{1-x^2}\cdot\left(1-x^2\right)}\cdot\arcsin x+\frac{x}{1-x^2}}
\end{align*}
lze ještě upravit takto:
\[
\frac{\arcsin x+x\cdot\sqrt{1-x^2}}{\sqrt{1-x^2}\cdot\left(1-x^2\right)} = \vysl{\frac{\arcsin x+x\cdot\sqrt{1-x^2}}{\sqrt{\left(1-x^2\right)^3}}}
\]

\hrule
\medskip
\textbf{\textcircled{10}}\quad $y=\tg\sqrt x+\underbrace{\tg\frac\pi4}_{\text{konstanta}}$
\[
y' = \frac{1}{\cos^2\sqrt x}\cdot\big(\underbrace{\sqrt x}_{x^{\frac12}}\big)'+0 = \frac{1}{\cos^2\sqrt x}\cdot\frac12\cdot x^{-\frac12} = \vysl{\frac{1}{2\sqrt x\cos^2\sqrt x}}
\]

\end{document}
